% partes/parte1/practicas-oficiales.tex
\chapter*{Prácticas oficiales de la unidad — Parte I}
\addcontentsline{toc}{chapter}{Prácticas oficiales de la unidad — Parte I}

Esta sección consolida, con nombre y horas oficiales del programa, las dos
prácticas de la Unidad I. Los insumos de cada práctica ya se desarrollaron
en las secciones de laboratorio de los capítulos 1--4; aquí se presenta la
guía formal de entrega.

\textbf{Metodología común a ambas prácticas.} En los cuatro capítulos, el
laboratorio usa \textbf{el mismo script} que el estudiante entrega en su
práctica: no hay una versión \enquote{de clase} y otra \enquote{de tarea}
distintas. El script no se modifica más allá del bloque de parámetros,
claramente delimitado al inicio de cada archivo (\texttt{matlab/capNN/*.m}).
La tarea del estudiante es correr ese mismo script para dos o tres casos
distintos (cambiando solo datos o parámetros), registrar los resultados que
imprime y grafica cada corrida, compararlos entre sí y redactar una
conclusión propia. Ningún entregable requiere que el estudiante escriba o
edite código MATLAB más allá de sustituir valores en el bloque de
parámetros.

\section*{Práctica 1 — El consumidor digital y su perfil (7 h)}

\begin{description}
  \item[Capítulos que la integran] 1 y 2.
  \item[Lugar de realización] Aula / laboratorio.
  \item[Objetivo] Analizar la evolución y el perfil del consumidor digital
  aplicando, sobre distintos casos generados solo cambiando parámetros, el
  modelo de difusión de Bass (Cap.~1) y una segmentación rígida y difusa
  (Cap.~2).
  \item[Entregable] Reporte único que integre: (a) el cuadro de los tres
  casos del ejercicio propuesto del Cap.~1 (\texttt{cap01\_difusion\_bass.m},
  Casos A/B/C de innovación vs. imitación) con las dos preguntas de
  comparación respondidas; (b) el cuadro de los tres casos del ejercicio
  propuesto del Cap.~2 (\texttt{cap02\_segmentacion.m}, Casos A/B/C de
  traslape y número de arquetipos) con las tres preguntas de comparación
  respondidas; (c) una conclusión propia de media página que relacione ambos
  laboratorios con el concepto de perfil como función del negocio (Cap.~2,
  §2.5).
  \item[Criterios de evaluación] Los dos scripts se ejecutan sin modificar
  nada fuera del bloque de parámetros (verificable por reproducibilidad,
  \texttt{rng(42)}); los tres casos de cada script están completos, no solo
  el caso por defecto; interpretación de negocio propia, no copiada de este
  libro.
\end{description}

\section*{Práctica 2 — Buscadores y plataformas en la gestión de datos del consumidor digital (8 h)}

\begin{description}
  \item[Capítulos que la integran] 3 y 4.
  \item[Lugar de realización] Aula / laboratorio.
  \item[Objetivo] Aplicar limpieza de datos, segmentación RFM, CLV y prueba
  A/B (Cap.~3) junto con análisis de series de tiempo y frecuencia de
  términos (Cap.~4) para fundamentar una decisión de negocio digital, en
  cada caso comparando resultados entre corridas del mismo script.
  \item[Entregable] Reporte único que integre: (a) el cuadro de los tres
  casos del ejercicio propuesto del Cap.~3
  (\texttt{cap03\_rfm\_clv\_ab.m}, Casos A/B/C de tamaño de muestra y nivel
  de significancia de la prueba A/B) con las tres preguntas de comparación
  respondidas; (b) el cuadro de los tres casos del ejercicio propuesto del
  Cap.~4 (\texttt{cap04\_series\_texto.m}, Casos A/B/C de amplitud
  estacional y ruido) junto con la clasificación de los diez términos más
  frecuentes del corpus de escucha social; (c) aplicación de la matriz de
  cuatro criterios del Cap.~4 a una plataforma digital real elegida por el
  equipo.
  \item[Criterios de evaluación] Rigor estadístico en la prueba A/B (uso
  correcto del valor $p$, no solo comparación de porcentajes, y explicación
  de por qué el mismo par de proporciones puede cambiar de decisión al
  cambiar el tamaño de muestra); aplicación correcta y honesta de la matriz
  de criterios (sin promoción implícita de ninguna plataforma); todas las
  cifras de mercado citadas en el entregable deben tener fuente y fecha de
  consulta.
\end{description}
